% !TeX root = ../../main.tex
\section{Bayes 滤波}\label{cw:sec:bayes}

\begin{lemma}[无记忆性]\label{cw:lem:markov}
在上述状态方程与噪声独立性条件下，给定 $X_{k-1}$ 后，$X_k$ 的条件分布不再依赖更早的状态与历史观测；给定 $X_k$ 后，$Y_k$ 的条件分布不再依赖过去的状态与观测。对 $k\geq1$ 有
\begin{equation}\label{cw:eq:2}
\begin{aligned}
&p_{X_k\mid X_{0:k-1},Y_{1:k-1}}(x_k\mid x_{0:k-1},y_{1:k-1})
=p_{X_k\mid X_{k-1}}(x_k\mid x_{k-1}),\\
&p_{Y_k\mid X_{0:k},Y_{1:k-1}}(y_k\mid x_{0:k},y_{1:k-1})
=p_{Y_k\mid X_k}(y_k\mid x_k).
\end{aligned}
\end{equation}
$k=1$ 时历史观测为空。第一式给出状态的一阶 Markov 性，第二式给出观测的条件独立性，分别用于下述预测与更新。

\end{lemma}
\begin{proof}
离散情形的证明及一般情形的推广说明见第~\ref{cw:app:markov-proof}~节。
\end{proof}

固定观测记录 $y_{1:k}=(y_1,\dots,y_k)^{\mathrm{T}}$，以 $p_{X_0}$ 初始化，$k=1$ 时不含历史观测条件。

\textbf{预测：}由上一时刻的后验与状态转移密度，求当前状态在收到 $y_k$ 前的条件密度。
\begin{equation}\label{cw:eq:3}
\begin{aligned}
&p_{X_k\mid Y_{1:k-1}}(x_k\mid y_{1:k-1})\\
= &\int_{\mathcal{X}}p_{X_k\mid X_{k-1}}(x_k\mid x_{k-1})p_{X_{k-1}\mid Y_{1:k-1}}(x_{k-1}\mid y_{1:k-1})\,\mathrm{d}v_X(x_{k-1}).
\end{aligned}
\end{equation}

\textbf{更新：}收到 $y_k$ 后，用观测似然对预测密度加权，再归一化，证明见定理~\ref{cw:thm:bayes}。
\begin{equation}\label{cw:eq:4}
\begin{aligned}
&p_{X_k\mid Y_{1:k}}(x_k\mid y_{1:k})\\
= &\frac{p_{Y_k\mid X_k}(y_k\mid x_k)p_{X_k\mid Y_{1:k-1}}(x_k\mid y_{1:k-1})}
{\displaystyle\int_{\mathcal{X}}p_{Y_k\mid X_k}(y_k\mid x)p_{X_k\mid Y_{1:k-1}}(x\mid y_{1:k-1})\,\mathrm{d}v_X(x)}.
\end{aligned}
\end{equation}
分母就是观测的预测密度 $p_{Y_k\mid Y_{1:k-1}}(y_k\mid y_{1:k-1})$，要求其有限且严格为正。更新所得后验作为下一步预测的输入，这两步构成精确的 Bayes 滤波递推。
